\documentclass{article}
\usepackage{mathnotes}

\title{Convexity}

\begin{document}

\section{Convexity}
\begin{definition}[Convex Function]\synonyms{convex}
Let $E \subseteq \reals^k$ be a \@{convex-set}. A function $\varphi : E \to \reals$ is \textbf{convex} if

\[ \varphi\big(\lambda \vec{x} + (1 - \lambda) \vec{y}\big) \leq \lambda \varphi(\vec{x}) + (1 - \lambda)\varphi(\vec{y}) \]

whenever $\vec{x}, \vec{y} \in E$ and $0 < \lambda < 1.$

In geometric terms, this means a \@{chord} joining any two points on the graph of $\varphi$ lies on or above the graph between them.
\end{definition}

\begin{definition}[Strictly Convex Function]
Starting with \@{convex-function}, $\varphi$ is \textbf{strictly convex} if the inequality is strict whenever $\vec{x} \neq \vec{y}.$
\end{definition}

\begin{definition}[Concave Function]\synonyms{concave}
A function $\varphi$ is said to be \textbf{concave} if $-\varphi$ is a \@{convex-function}.
\end{definition}

\begin{theorem}[Adding an Affine Function Preserves Convexity]
Let $E \subseteq \reals^k$ be a \@{convex-set}, let $\varphi : E \to \reals,$ and let $\ell : E \to \reals$ be \@{affine}. If $\varphi$ is \@{convex}, then $\varphi + \ell$ and $\varphi - \ell$ are convex. If $\varphi$ is \@{concave}, then $\varphi + \ell$ and $\varphi - \ell$ are concave. A linear function is affine, so this covers adding or subtracting a linear function.

TODO: prove this. $\ell$ satisfies the convexity inequality with equality, so adding it to both sides of the inequality for $\varphi$ leaves the inequality intact; subtraction is the same argument applied to $-\ell,$ which is also affine.
\end{theorem}

\begin{theorem}[Affine Precomposition Preserves Convexity]
Let $D \subseteq \reals^k$ and $E \subseteq \reals^m$ be \@{convex-set}s, let $A : D \to E$ be \@{affine}, and let $\varphi : E \to \reals.$ If $\varphi$ is \@{convex}, then $\varphi \circ A$ is convex on $D.$ If $\varphi$ is \@{concave}, then $\varphi \circ A$ is concave on $D.$

TODO: prove this. $A$ carries a \@{convex-combination} of two points to the same convex combination of their images, so the inequality for $\varphi$ applies to $A(\vec{x})$ and $A(\vec{y})$ directly.
\end{theorem}

\begin{theorem}[Second Derivative Test for Convexity]\label{second-derivative-test-for-convexity}
Let $\varphi$ be twice \@{differentiable} on an \@{open} \@{interval} $(a,b).$ Then
$\varphi$ is \@{convex} on $(a,b)$ if and only if $\varphi''(x) \geq 0$ for all
$x \in (a,b),$ and \@{concave} if and only if $\varphi''(x) \leq 0$ for all
$x \in (a,b).$
\end{theorem}

\begin{definition}[Convex Combination]
A \textbf{convex combination} is a \@{linear-combination} of \@{points} where all \@{coefficients} are non-negative and sum to 1.
\end{definition}

\begin{note}
Every \@{convex-combination} of two \@{points} lies on the \@{line-segment} between the \@{points}.
\end{note}

\end{document}
